\section{Implementation}
\label{sec:implementation}

In this section, we describe the implementation of the toolset.
The SAT encoding we use is introduced in \cite{listenmaa_claessen2015}, along with a working CG compiler.
Since we want to test the rules in a way that doesn't depend on data, we operate on a \emph{symbolic sentence}: a word sequence which we can modify according to the properties we want to test.

\paragraph{SAT encoding of CG rules}

A single CG rule is a template which creates a concrete clause for each instance where it is applied.
%Each potential analysis is represented as a variable, and the rules are translated into clauses that operate on those variables.
Table~\ref{table:vars} shows the rule \texttt{REMOVE v IF (-1 det)} applied to the text \emph{the bear}.
The analyses of the word forms are represented as variables, and
the default rule ensures that each word has at least one analysis.
%This makes \texttt{the\_det} true by default, therefore the implication must hold, and \texttt{bear\_v\_pl} will be removed.
For more details about the encoding, see \cite{listenmaa_claessen2015}.

% table:vars moved to main file for positioning


\paragraph{Word-internal constraints}

The symbolic sentence consists of symbolic words.
In the beginning, a symbolic word is a list of all possible tag combinations, and we don't know which ones are true.
%These tag combinations may be underspecified for lemma or other features---similarly, a CG rule may target any noun, a plural noun or the noun \emph{bear}.
Given a large morphological lexicon, we can restrict the individual words to contain only realistic ambiguities. 
For instance, a Spanish morphological lexicon contains the entries \texttt{casa:casar<verb><sg><p3>} and \texttt{casa:casa<noun><sg>}, hence we know that a confusion between a 3rd person singular verb and a singular noun is possible. 
%The lexicon is not likely to contain a wordform that can be analysed both as an adjective and punctuation, so that kind of ambiguous word is never created.

%However, since our method is corpus-free, we don't add any external restrictions as to what words can follow each other.

\paragraph{Symbolic sentences}

We apply a concrete rule to a symbolic sentence
we get back a symbolic sentence: it's the same table except that 

we create a new variable for w2<verb> (target)a
 with the target changed into a new variable

definition of w2'<verb> = *nice formula*


Figure~\ref{fig:example} shows a fragment of of a two-word symbolic sentence, before and after applying the rule \texttt{REMOVE verb IF (-1 det)}\footnote{For simplicity, we don't show the word-internal constraints that come from the lexicon, but only the rule application.}.
As with real analyses, each tag combination is given a variable, and applying a rule means creating a concrete clause for each instance where a determiner is found preceding a verb.
In a symbolic sentence, every word has potentially every tag combination, and applying the rules will shape the sentence into something that respects the constraints of the language.

\begin{figure}[h]
\includegraphics[width=0.42\textwidth]{symbolic_sentence.png}
\caption{Symbolic sentence beginning to take shape \todo{make this less confusing?}}
\label{fig:example}
\end{figure}


The example rule states that the sentence may not have a verb after a determiner.
The variable \texttt{w2<verb>} is the situation of the second word's verb status before applying the rule. We create a new variable \texttt{w2'<verb>}, whose value depends on whether the condition holds (if \texttt{w1<det>} is true); whether it was true on the last round (if \texttt{w2<verb>} is already removed by another rule, \texttt{w2'<verb>} cannot reappear); and whether it is the only analysis left (not depicted in the picture).
All the analyses not targeted by this rule remain the same.


% \begin{figure}[]
% \includegraphics[width=0.42\textwidth]{symbolic_sentence.png}
% \caption{Symbolic sentence beginning to take shape \todo{make this less confusing?}}
% \label{fig:example}
% \end{figure}

%The applications of these rules shape our symbolic sentence into containing only things that can survive the rules. 

The effect of the old rule applications is retained when a new rule is added.
This means that we can ask a question ``after applying rules $0-i$, is it possible for $i+1$ to apply?'', and the answer is reliable; we cannot find a counterexample by just having a bit more data, because we start from literally all possibilities, and each rule restricts the set of possible sentences.
%If rule \texttt{r100} would require the first word to be a noun, and causes a conflict because it is not a noun, we know that it is because there is no way for w1 to be a noun.

%Note that this method is not suitable for language generation: rules in any realistic CG are tailored for real life ambiguities, whereas our symbolic sentence starts with literally all possibilities of tag combinations.

\paragraph{What kind of questions we can ask?}

The example so far has been of type ``Can rule Z apply after A-Y?''
In case a conflict is detected, we can dig into it and refine where the problem is: e.g. is the condition not met, or is it the case that the target is already removed or the only remaining analysis.
We can find what other rule or rules prevent others from applying, or if the conflict is rule-internal, such as nonexisting tag set or contradicting requirements of a context word (e.g. a word must be unambiguously two different POS).


In addition to analysing a whole grammar, 
we can construct all kinds of symbolic sentences to test out individual rules. 
We can set the length, restrict individual words (e.g. ``3rd word must be a noun''), require that it triggers some rule but not other. 
This functionality can aid the grammar writer in the process, to see if they have missed a case or defined the tagset correctly.






